% Generated by 10_final_submission_estimates.R -- do not edit by hand.
\begin{table}[htbp]
\centering
\caption{Summary statistics: boundary-level characteristics}
\label{tab:summary_stats}
\begin{threeparttable}
\small
\begin{tabular}{lcccccc}
\toprule
Variable & Mean & Std.\ dev. & Min & Median & Max \\
\midrule
\addlinespace[0.1cm]
\multicolumn{6}{l}{\textit{Panel A: Sample composition}} \\
\addlinespace[0.1cm]
Incidents within 2,000 m of the boundary & 1,368 & 1,091 & 336 & 867 & 4,082 \\
Dates on which the boundary is active & 70 & 1 & 67 & 70 & 72 \\
Estimation cells (date $\times$ bin) & 5,600 & 109 & 5,360 & 5,600 & 5,760 \\
Cells with positive kernel weight & 1,428 & 364 & 980 & 1,340 & 2,240 \\
\addlinespace[0.2cm]
\midrule
\addlinespace[0.1cm]
\multicolumn{6}{l}{\textit{Panel B: Bandwidth}} \\
\addlinespace[0.1cm]
MSE-optimal bandwidth (m) & 504 & 129 & 335 & 481 & 797 \\
\addlinespace[0.2cm]
\midrule
\addlinespace[0.1cm]
\multicolumn{6}{l}{\textit{Panel C: Boundary-specific treatment effects}} \\
\addlinespace[0.1cm]
Point estimate ($\hat{\tau}_b$) & 0.0181 & 0.1009 & -0.1352 & -0.0002 & 0.2279 \\
Standard error & 0.0646 & 0.0275 & 0.0093 & 0.0584 & 0.1127 \\
Absolute effect size & 0.0710 & 0.0716 & 0.0002 & 0.0500 & 0.2279 \\
\bottomrule
\end{tabular}
\begin{tablenotes}[flushleft]
\small
\item \textit{Notes:} Statistics across the 15 boundaries that generate identifying variation, unweighted. Incidents are counted within 2{,}000 m of the assigned boundary and total 20,527. Estimation cells are boundary $\times$ date $\times$ 50-metre-bin units of the balanced zero-filled panel and total 84,000; 21,414 of them carry positive kernel weight at the MSE-optimal bandwidth. A boundary is active only on dates when one of its two polygons is rationed, which is why the date counts are far below the 325-day study window.
\end{tablenotes}
\end{threeparttable}
\end{table}
